% !TeX root = main.tex


\section{Complexity Analysis}

The computational complexity of gPCD is determined primarily by
broad-phase occupancy generation and narrow-phase collision
evaluation. The complexity results summarized in
Table~\ref{tab:ComplexityTable} are derived below.


\begin{table}[h!]
\centering
\caption{Summary of gPCD computational and storage complexities.}
\label{tab:complexity_summary}
\begin{tabular}{lcc}
\hline
Component & Complexity & Type \\
\hline
Corner Evaluation & $O(1)$ & Computational \\
Cell Determination & $O(1)$ & Computational \\
Broad-Phase Occupancy Generation & $O(N)$ & Computational \\
Narrow-Phase Evaluation & $O(K)$ & Computational \\
Particle Storage & $O(N)$ & Storage \\
Cell-Array Storage & $O(S^3)$ & Storage \\
\hline
\end{tabular}
\end{table}



Each particle performs a fixed number of occupancy evaluations
independent of the total particle count. Consequently, broad-phase
occupancy generation scales linearly with the number of particles,

\begin{equation}
T_{\mathrm{occ}}(N)=O(N).
\end{equation}

Because occupancy lists are accessed directly through flattened cell
indices, cell lookup requires constant-time addressing. Under the fixed
occupancy-width assumptions used by the framework, insertion and
retrieval therefore require bounded work per operation.

Each source particle traverses the occupancy lists associated with its
corner cells. Let $K$ denote the total number of interaction
evaluations performed during narrow-phase execution. Since each
interaction requires bounded work,

\begin{equation}
T_{\mathrm{np}}(K)=O(K).
\end{equation}

Consequently, narrow-phase complexity is governed by the number of
candidate interactions (K) represented within the potentially
colliding set.

For comparison, a naive all-pairs collision search that evaluates each
unordered particle pair once requires

\begin{equation}
K_{\mathrm{all}} =\frac{N(N-1)}{2}=O(N^2)
\label{niaeve}
\end{equation}

pair evaluations.

The source-owned narrow phase used by gPCD does not eliminate the
theoretical worst-case quadratic bound associated with pairwise
collision testing. If every particle is a candidate for every other
particle, then \(K=O(N^2)\). However, candidate evaluation remains
source-owned throughout the narrow phase, avoiding construction of a
globally shared collision-pair structure and reducing synchronization
requirements during massively parallel execution.

Combining broad-phase occupancy generation and narrow-phase evaluation
gives

\begin{equation}
T_{\mathrm{gPCD}} = O(N+K),
\end{equation}

where \(K\) is the number of candidate interactions represented within
the potentially colliding set.

Particle storage scales linearly with particle count,

\begin{equation}
M_{\mathrm{particle}}=O(N).
\end{equation}

For a cubic cell array of side length (S) and occupancy-list width
(W),

\begin{equation}
M_{\mathrm{cell}}=O(S^3W).
\end{equation}

The total storage requirement is therefore

\begin{equation}
M_{\mathrm{gPCD}}=O(N+S^3W),
\end{equation}

which reduces to $(O(N+S^3))$ for fixed occupancy width.



